\documentclass[11pt,a4paper]{article}

% ------------------------------------------------------------------ Encodage
\usepackage[utf8]{inputenc}
\usepackage[T1]{fontenc}
\usepackage{lmodern}
\usepackage[french]{babel}

% --------------------------------------------------------------------- Mise en page
\usepackage[a4paper,margin=2.4cm]{geometry}
\usepackage{xcolor}
\usepackage{graphicx}
\usepackage{booktabs}
\usepackage{tabularx}
\usepackage{array}
\usepackage{longtable}
\usepackage{enumitem}
\usepackage{fancyhdr}
\usepackage{titlesec}
\usepackage{listings}
\usepackage[hidelinks]{hyperref}

% ---------------------------------------------------------------------- Couleurs
\definecolor{primary}{HTML}{1F4E79}
\definecolor{accent}{HTML}{2E75B6}
\definecolor{ok}{HTML}{1E7D32}
\definecolor{warn}{HTML}{C77700}
\definecolor{danger}{HTML}{B3261E}
\definecolor{codebg}{HTML}{F4F6F8}
\definecolor{rule}{HTML}{C9D6E2}

\hypersetup{
  colorlinks=true,
  linkcolor=primary,
  urlcolor=accent,
  pdftitle={Rapport de deploiement - GANDAL-API},
  pdfauthor={Delmat237}
}

% ------------------------------------------------------------------ En-tete / pied
\pagestyle{fancy}
\fancyhf{}
\renewcommand{\headrulewidth}{0.4pt}
\renewcommand{\footrulewidth}{0.4pt}
\fancyhead[L]{\small\textcolor{primary}{\textbf{GANDAL-API}} -- Rapport de déploiement}
\fancyhead[R]{\small 11 juin 2026}
\fancyfoot[L]{\small\textit{Document confidentiel}}
\fancyfoot[R]{\small Page \thepage}

% ----------------------------------------------------------------- Titres de section
\titleformat{\section}
  {\color{primary}\Large\bfseries}{\thesection}{0.6em}{}
\titleformat{\subsection}
  {\color{accent}\large\bfseries}{\thesubsection}{0.6em}{}
\titlespacing*{\section}{0pt}{1.4em}{0.7em}

% --------------------------------------------------------------------- Listings
\lstdefinestyle{cmd}{
  basicstyle=\ttfamily\footnotesize,
  backgroundcolor=\color{codebg},
  frame=single,
  rulecolor=\color{rule},
  breaklines=true,
  columns=fullflexible,
  keepspaces=true,
  showstringspaces=false,
  xleftmargin=0.5em,
  framesep=6pt,
  aboveskip=1em,
  belowskip=1em
}
\lstset{style=cmd}

% --------------------------------------------------------------------- Macros
\newcommand{\code}[1]{\texttt{\small #1}}
\newcommand{\ok}{\textcolor{ok}{\textbf{OK}}}
\newcommand{\statusbadge}[2]{\colorbox{#1}{\textcolor{white}{\textbf{\,#2\,}}}}

\setlist{topsep=2pt,itemsep=1pt}
\renewcommand{\arraystretch}{1.25}

% =====================================================================
\begin{document}

% --------------------------------------------------------------- Page de titre
\begin{titlepage}
\centering
\vspace*{1.5cm}

{\color{rule}\rule{\linewidth}{1.2pt}}\\[0.6cm]
{\color{primary}\Huge\bfseries Rapport de déploiement}\\[0.45cm]
{\LARGE GANDAL-API \textnormal{(\code{dc-backend})}}\\[0.35cm]
{\large\itshape API de gestion de VMs pédagogiques, requêtes et publications}\\[0.5cm]
{\color{rule}\rule{\linewidth}{1.2pt}}\\[1.2cm]

\statusbadge{ok}{EN PRODUCTION -- OPÉRATIONNEL ET VALIDÉ}\\[1.6cm]

\begin{tabular}{@{}r l@{}}
\textbf{Projet} & dc-backend (\href{https://github.com/Delmat237/GANDAL-API}{github.com/Delmat237/GANDAL-API})\\
\textbf{Version} & commit \code{1a48612} -- branche \code{master}\\
\textbf{Serveur cible} & \code{omega-test-3001} -- \code{10.50.30.102}\\
\textbf{Type} & Conteneurisé (Docker Compose), en isolation\\
\textbf{Auteur} & Cellule DevOps\\
\textbf{Date} & 11 juin 2026\\
\end{tabular}

\vfill
{\small\color{primary}Document technique d'ingénierie -- diffusion restreinte}
\end{titlepage}

% --------------------------------------------------------------- Sommaire
\tableofcontents
\thispagestyle{fancy}
\newpage

% =====================================================================
\section{Résumé exécutif}

L'API GANDAL a été déployée avec succès sur le serveur \code{10.50.30.102} sous forme
de \emph{stack} Docker Compose isolée, \textbf{sans aucun impact} sur les applications
tierces déjà hébergées sur la machine (\emph{stacks} \code{iwm-*}, \code{erp-cashier-*},
\code{verifid}).

La solution expose l'API selon deux canaux complémentaires :
\begin{enumerate}
  \item \textbf{Reverse-proxy mutualisé} (nginx \code{iwm-gateway}, port 80) -- accès
        nominal de production via le sous-domaine \code{gandal.10.50.30.102.nip.io} ;
  \item \textbf{Port hôte direct} \code{8000} -- accès de diagnostic.
\end{enumerate}

L'ensemble de la chaîne a été vérifié de bout en bout : démarrage des conteneurs,
exécution des migrations de base de données, points de santé applicatifs, routage du
reverse-proxy, création du compte d'administration et authentification.
\textbf{Tous les tests de recette sont au vert.}

% =====================================================================
\section{Contexte et objectifs}

\begin{tabularx}{\linewidth}{@{}X c@{}}
\toprule
\textbf{Objectif} & \textbf{Atteint} \\
\midrule
Déployer l'API sur \code{root@10.50.30.102} & \ok \\
Ne pas perturber les services existants du serveur & \ok \\
Exposer l'API derrière le reverse-proxy nginx existant & \ok \\
Sécuriser les secrets de production (JWT, mot de passe admin, base) & \ok \\
Disposer d'un compte administrateur fonctionnel & \ok \\
\bottomrule
\end{tabularx}

% =====================================================================
\section{Environnement cible}

\begin{tabularx}{\linewidth}{@{}l X@{}}
\toprule
\textbf{Élément} & \textbf{Valeur} \\
\midrule
Hôte & \code{omega-test-3001} \\
Adresse IP & \code{10.50.30.102} \\
Système & Debian GNU/Linux 12 (bookworm) \\
Noyau & \code{6.1.0-49-cloud-amd64} \\
Moteur de conteneurs & Docker \code{29.5.3} \\
Orchestration & Docker Compose \code{v5.1.4} (plugin) \\
Stockage racine & 20 Go ($\approx$ 87 \% utilisés après déploiement) \\
Répertoire applicatif & \code{/gandal-dev} \\
\bottomrule
\end{tabularx}

\medskip
\noindent\textbf{Note environnement partagé.} Le serveur héberge plusieurs \emph{stacks}
indépendantes (\code{iwm-app}, \code{iwm-postgres}, \code{iwm-kafka},
\code{iwm-elasticsearch}, \code{iwm-redis}, \code{erp-cashier-backend},
\code{erp-cashier-frontend}, \code{verifid-backend}, reverse-proxy \code{iwm-gateway}).
Le déploiement a été conduit en stricte isolation.

\newpage{}
% =====================================================================
\section{Architecture de la solution déployée}

\begin{lstlisting}
                         Internet / LAN
                               |
              http://gandal.10.50.30.102.nip.io   (port 80)
                               |
                 +-------------v--------------+
                 |   nginx  "iwm-gateway"     |  (reverse-proxy mutualise)
                 |  routage par sous-domaine  |
                 |  map $host $svc {          |
                 |    ~^gandal\. -> gandal-api |  <- entree ajoutee
                 |    ~^api\.    -> iwm-app     |
                 |    ~^caisse\. -> ...         |
                 |  }                          |
                 +-------------+--------------+
                               | reseau Docker "iwm-vm-network"
                               |
            +------------------v-------------------+
            |  Stack Compose  "gandal"             |
            |                                       |
            |  +---------------+   +--------------+ |
   :8000 ---+->|  gandal-api   |-->|   gandal-db  | |
  (acces    |  |  FastAPI      |   | postgres:17  | |
  direct)   |  |  172.18.0.10  |   |  (healthy)   | |
            |  +---------------+   +------+-------+ |
            |      reseau "gandal-net"    |        |
            |                      volume gandal_pgdata
            +---------------------------------------+
\end{lstlisting}

\subsection*{Inventaire des conteneurs déployés}

\begin{tabularx}{\linewidth}{@{}l l c X c@{}}
\toprule
\textbf{Conteneur} & \textbf{Image} & \textbf{Taille} & \textbf{Réseaux} & \textbf{Statut} \\
\midrule
\code{gandal-api} & build local & 311 Mo & \code{gandal-net}, \code{iwm-vm-network} & Up \\
\code{gandal-db} & \code{postgres:17-alpine} & 400 Mo & \code{gandal-net} & Up (healthy) \\
\bottomrule
\end{tabularx}

\medskip
\noindent\textbf{Persistance} : volume Docker nommé \code{gandal\_pgdata} (données
PostgreSQL).

% =====================================================================
\section{Procédure de déploiement réalisée}

\begin{enumerate}[label=\textbf{\arabic*.}]
  \item \textbf{Audit préalable du serveur} -- vérification de l'accès SSH, de la
        présence de Docker / Compose, des ports hôtes disponibles (80 occupé par le
        gateway ; 8000 et 5433 libres) et de l'absence d'un déploiement antérieur.
  \item \textbf{Analyse du reverse-proxy existant} -- identification du mécanisme de
        routage par sous-domaine (\code{map \$host \$svc}) et de la résolution DNS
        Docker sur le réseau \code{iwm-vm-network}.
  \item \textbf{Transfert du code} -- synchronisation du dépôt local vers
        \code{/gandal-dev} via \code{rsync} (exclusions : \code{.git/}, \code{.venv/},
        caches, \code{docs/}, \code{.env}).
  \item \textbf{Génération des secrets de production} -- JWT, clé secrète SuperAdmin,
        mot de passe PostgreSQL et mot de passe administrateur (valeurs fortes
        aléatoires).
  \item \textbf{Définition de la \emph{stack} de production} -- fichier
        \code{docker-compose.prod.yml} dédié (projet \code{gandal}, conteneurs nommés,
        raccordement au réseau du gateway, volume persistant, sondes de santé).
  \item \textbf{Construction et démarrage} -- \code{docker compose up -d -{}-build} ;
        les migrations Alembic (\code{alembic upgrade head}) s'exécutent automatiquement
        au démarrage du conteneur API.
  \item \textbf{Intégration au reverse-proxy} -- ajout d'une entrée de routage
        \code{gandal.} dans \code{/root/iwm/ops/nginx/gateway.conf} (avec sauvegarde
        préalable), validation \code{nginx -t} puis rechargement à chaud
        \code{nginx -s reload}.
  \item \textbf{Bootstrap applicatif} -- création du compte SuperAdmin via
        \code{POST /api/v1/auth/register-admin}.
  \item \textbf{Recette} -- validation des points de santé, du routage et de
        l'authentification.
\end{enumerate}

% =====================================================================
\section{Configuration}

\subsection{Stack Docker Compose}

Caractéristiques notables de \code{/gandal-dev/docker-compose.prod.yml} :
\begin{itemize}
  \item \code{name: gandal} -- projet Compose isolé ;
  \item conteneurs nommés explicitement (\code{gandal-api}, \code{gandal-db}) pour la
        résolution DNS par le reverse-proxy ;
  \item \code{restart: unless-stopped} -- redémarrage automatique des services ;
  \item sonde de santé PostgreSQL (\code{pg\_isready}) ; l'API ne démarre qu'une fois la
        base \code{healthy} (\code{depends\_on: condition: service\_healthy}) ;
  \item raccordement de \code{gandal-api} au réseau externe \code{iwm-vm-network} du
        gateway \textbf{et} à un réseau privé \code{gandal-net} ;
  \item publication de l'API sur \code{0.0.0.0:8000} (accès direct de diagnostic).
\end{itemize}

\subsection{Variables d'environnement (\code{.env})}

\begin{tabularx}{\linewidth}{@{}l X@{}}
\toprule
\textbf{Variable} & \textbf{Valeur de production} \\
\midrule
\code{ENVIRONMENT} & \code{production} \\
\code{DATABASE\_URL} & \code{postgresql://dc:***@db:5432/dc\_backend} \\
\code{JWT\_SECRET} & secret fort généré \textit{(masqué)} \\
\code{SUPERADMIN\_USERNAME} & \code{admin} \\
\code{SUPERADMIN\_EMAIL} & \code{admin@enspy-gi.cm} \\
\code{SUPERADMIN\_PASSWORD} & secret fort généré \textit{(transmis séparément)} \\
\code{SUPERADMIN\_SECRET\_KEY} & secret fort généré \textit{(masqué)} \\
\code{PROXMOX\_ENABLED} & \code{false} (mode simulation) \\
\code{EMAIL\_BACKEND} & \code{log} \\
\code{CORS\_ORIGINS} & origines \code{localhost} de développement \\
\bottomrule
\end{tabularx}

% =====================================================================
\section{Intégration au reverse-proxy nginx}

Une seule ligne a été ajoutée à la table de routage du gateway mutualisé,
\code{/root/iwm/ops/nginx/gateway.conf} :

\begin{lstlisting}[language=]
map $host $svc {
      default               "";
      "~^gandal\."          "gandal-api:8000";   # <- ajout GANDAL
      "~^api-caisse\."      "erp-cashier-backend:8081";
      "~^api\."             "iwm-app:8080";
      "~^verifid\."         "verifid-backend:8080";
      "~^caisse\."          "erp-cashier-frontend:3000";
}
\end{lstlisting}

\begin{itemize}
  \item \textbf{Sauvegarde} : une copie horodatée \code{gateway.conf.bak.<timestamp>} a
        été créée avant modification ;
  \item \textbf{Validation} : \code{nginx -t} $\rightarrow$ \emph{syntax is ok / test is
        successful} ;
  \item \textbf{Application} : rechargement à chaud (\code{nginx -s reload}), \textbf{sans
        interruption} des autres services ;
  \item \textbf{Non-régression} : un \code{Host} inconnu continue de renvoyer
        \code{404 Unknown host}, confirmant que les routes existantes restent intactes.
\end{itemize}

% =====================================================================
\section{Sécurité et gestion des secrets}

\begin{itemize}
  \item Tous les secrets de production ont été \textbf{régénérés} (aucune valeur
        d'exemple \code{changeme123} / \code{change-me-*} conservée) ;
  \item les secrets résident dans \code{/gandal-dev/.env} sur le serveur ; ce fichier
        \textbf{n'est pas versionné} (exclu par \code{.gitignore} et du transfert
        \code{rsync}) ;
  \item le mot de passe administrateur a été communiqué \textbf{hors de ce document} ;
  \item la base de données PostgreSQL n'est \textbf{pas exposée} sur l'hôte (port interne
        au réseau Docker uniquement) ;
  \item le compte d'administration est protégé par une clé secrète
        (\code{SUPERADMIN\_SECRET\_KEY}) exigée à la création
        (\code{secrets.compare\_digest}).
\end{itemize}

% =====================================================================
\section{Recette -- tests de validation}

\begin{longtable}{@{}c >{\raggedright\arraybackslash}p{0.42\linewidth} >{\raggedright\arraybackslash}p{0.30\linewidth} c@{}}
\toprule
\textbf{\#} & \textbf{Test} & \textbf{Résultat attendu} & \textbf{Statut} \\
\midrule
\endhead
1 & Démarrage base (\code{gandal-db}) & \code{Up (healthy)} & \ok \\
2 & Migrations Alembic & révisions \code{0001 -> 0003} & \ok \\
3 & Démarrage API & \code{Application startup complete} & \ok \\
4 & Santé -- \code{GET /health} & \code{\{"status":"ok"\}} & \ok \\
5 & Disponibilité DB -- \code{GET /ready} & \code{database: true} & \ok \\
6 & Routage gateway (Host \code{gandal.*}) & \code{\{"status":"ok"\}} & \ok \\
7 & Non-régression (Host inconnu) & \code{404 Unknown host} & \ok \\
8 & Documentation -- \code{GET /docs} & \code{HTTP 200} (Swagger) & \ok \\
9 & Bootstrap admin -- \code{register-admin} & \code{HTTP 201} & \ok \\
10 & Authentification -- \code{login} & \code{HTTP 200} + jeton JWT & \ok \\
11 & Accès port direct \code{:8000/docs} & \code{HTTP 200} & \ok \\
12 & Accès via nip.io \code{/docs} & \code{HTTP 200} & \ok \\
\bottomrule
\end{longtable}

\newpage{}
% =====================================================================
\section{Branchement de l'hyperviseur Proxmox (intégration réelle)}

Au-delà du \emph{mode simulation} utilisé pour la mise en production initiale, le module
d'intégration Proxmox de GANDAL a été \textbf{raccordé et validé contre un cluster
Proxmox réel}. Cette section documente le branchement, les ajustements de code rendus
nécessaires par le comportement réel de l'API, et le test de clonage de bout en bout.

\subsection{Topologie du cluster cible}

\begin{tabularx}{\linewidth}{@{}l X@{}}
\toprule
\textbf{Élément} & \textbf{Valeur} \\
\midrule
Hyperviseur & Proxmox VE \code{9.2.2} (cluster à 7 nœuds) \\
Nœud d'API / hôte & \code{emilia} -- \code{192.168.123.100:8006} \\
Autres nœuds & \code{ram} (.101), \code{rem} (.102), \code{BLADE} (.103),
\code{gandal} (.104), \code{GENESIS} (.105), \code{Bris} (.106) \\
Nœud de provisionnement & \code{emilia} (templates locaux, clonage local fiable) \\
Templates disponibles & \code{2299} (\code{omega-cloud-template}, cloud-init),
\code{9001} (\code{omega-template-base}) \\
Stockages & \code{stockage.ceph} (RBD partagé, 3{,}4 To), \code{local-lvm} (337 Go) \\
Pont réseau (emilia) & \code{vmbr2} \\
\bottomrule
\end{tabularx}

\medskip
\noindent\textbf{Choix du nœud.} Le proxy inter-nœuds \emph{depuis} \code{emilia} renvoie
une erreur \code{596 (Broken pipe)} ; le nœud \code{emilia} a donc été retenu comme cible
de provisionnement, les templates y résidant localement (clonage sans proxy). Le cluster
comporte par ailleurs un nœud nommé \code{gandal} (\code{192.168.123.104}), utilisable
ultérieurement en pointant directement son IP et en s'appuyant sur le stockage \code{ceph}
partagé.

\subsection{Génération du jeton API et configuration}

Un jeton d'API dédié a été généré sur \code{192.168.123.100} pour l'utilisateur
\code{root@pam} :
\begin{itemize}
  \item identifiant de jeton : \code{root@pam!gandal} ;
  \item \code{privsep=0} -- le jeton hérite des privilèges complets (clone, configuration,
        démarrage, suppression) ;
  \item le secret n'est affiché qu'une seule fois à la création ; il est stocké dans
        \code{/.env} \textbf{non versionné} (exclu par \code{.gitignore}).
\end{itemize}

\noindent Variables \code{.env} correspondantes (secret masqué) :

\begin{tabularx}{\linewidth}{@{}l X@{}}
\toprule
\textbf{Variable} & \textbf{Valeur} \\
\midrule
\code{PROXMOX\_ENABLED} & \code{true} \\
\code{PROXMOX\_HOST} & \code{192.168.123.100} \\
\code{PROXMOX\_USER} & \code{root@pam} \\
\code{PROXMOX\_TOKEN\_ID} & \code{gandal} \\
\code{PROXMOX\_TOKEN\_SECRET} & \textit{(secret masqué -- transmis séparément)} \\
\code{PROXMOX\_DEFAULT\_NODE} & \code{emilia} \\
\code{PROXMOX\_VM\_STORAGE} & \code{local-lvm} \\
\code{PROXMOX\_NET0\_TEMPLATE} & \code{virtio,bridge=vmbr2} \\
\code{PROXMOX\_VMID\_RANGE\_START} / \code{\_END} & \code{2400} / \code{2499} \\
\code{PROXMOX\_CLONE\_FULL} & \code{false} (clone lié) \\
\code{PROXMOX\_CLONE\_TIMEOUT} & \code{300} (secondes) \\
\bottomrule
\end{tabularx}

\subsection{Plage de VMID dédiée}

Une plage \textbf{\code{2400--2499}} a été réservée à GANDAL et vérifiée \textbf{entièrement
libre}. Elle évite toute collision avec les VMs existantes du cluster (gabarits
\code{2299}, flotte \code{omega} en \code{2300+} et \code{3000+}, services en \code{100+}).
La méthode \code{get\_next\_vmid()} alloue désormais le premier identifiant libre
\emph{dans cette plage} (au lieu de \code{cluster/nextid} global).

\subsection{Ajustements apportés au client Proxmox}

Le comportement réel de l'API a révélé plusieurs écarts par rapport à l'implémentation
initiale, tous corrigés dans \code{app/infrastructure/proxmox/proxmox\_client.py} :

\begin{longtable}{@{}p{0.30\linewidth} p{0.62\linewidth}@{}}
\toprule
\textbf{Ajustement} & \textbf{Motif / comportement réel} \\
\midrule
\endhead
Attente de fin de tâche de clonage & Le clonage est asynchrone (UPID) ; la VM reste
\emph{verrouillée} (\code{lock=clone}) pendant la copie. Le code attend désormais la fin
de la tâche avant de configurer. \\
Clone lié par défaut (\code{full=0}) & Un clone \emph{complet} de \code{ceph} vers
\code{local-lvm} (disque 20 Go) dépassait 5 minutes ; le clone lié s'effectue en quelques
secondes (disque conservé sur le \code{ceph} partagé). \\
Clé SSH URL-encodée & Sans encodage, l'API rejette le champ (\code{invalid urlencoded
string}). La clé est désormais encodée via \code{urllib.parse.quote}. \\
Configuration via \code{PUT} & Sur une VM fraîchement clonée, \code{POST .../config}
plaçait les paramètres matériels (RAM, vCPU, réseau) en \emph{pending} sans les appliquer ;
le \code{PUT} synchrone les applique immédiatement. \\
Suppression robuste & \code{delete\_vm} attend l'arrêt effectif de la VM (l'arrêt étant
asynchrone) puis supprime avec \code{purge}, évitant l'erreur \code{VM is running}. \\
\bottomrule
\end{longtable}

\subsection{Test de clonage réel -- bout en bout}

Un cycle de vie complet a été exécuté \textbf{via le client applicatif lui-même} (lecture
du \code{.env}, jeton réel), puis la VM de test a été détruite :

\begin{longtable}{@{}c >{\raggedright\arraybackslash}p{0.52\linewidth} >{\raggedright\arraybackslash}p{0.20\linewidth} c@{}}
\toprule
\textbf{\#} & \textbf{Étape} & \textbf{Résultat} & \textbf{Statut} \\
\midrule
\endhead
1 & Connexion au cluster (jeton) & Proxmox VE \code{9.2.2} & \ok \\
2 & Allocation VMID (\code{get\_next\_vmid}) & \code{2400} (dans la plage) & \ok \\
3 & Clone lié du template \code{2299} & disque \code{vm-2400-disk-0} & \ok \\
4 & Configuration appliquée (\code{PUT}) & \code{2048} Mo, \code{2} vCPU, \code{vmbr2},
clé SSH & \ok \\
5 & Démarrage (\code{start\_vm}) & \code{status=running} & \ok \\
6 & Suppression (\code{delete\_vm}) & VM purgée, cluster propre & \ok \\
\bottomrule
\end{longtable}

\noindent\textbf{Durée totale clone + configuration : $\approx$ 15 secondes.} L'ensemble de
la suite de tests applicative (\code{pytest}) reste au vert (\textbf{102 tests}).

\subsection{Activation en production}

L'intégration a été développée et validée depuis le réseau de laboratoire
(\code{192.168.123.0/24}). Pour l'activer sur le serveur de production
(\code{10.50.30.102}), il suffit de :
\begin{enumerate}
  \item s'assurer que l'hôte de production atteint \code{192.168.123.100:8006} (routage /
        pare-feu) ;
  \item reporter les variables \code{PROXMOX\_*} ci-dessus dans le \code{.env} serveur ;
  \item basculer \code{PROXMOX\_ENABLED=true} et \code{PROXMOX\_SIMULATION\_MODE=false}.
\end{enumerate}
\noindent Aucune autre modification de code n'est requise.

% =====================================================================
\section{Accès à l'application}

\begin{tabularx}{\linewidth}{@{}l X l@{}}
\toprule
\textbf{Canal} & \textbf{URL} & \textbf{Usage} \\
\midrule
Reverse-proxy (nominal) & \url{http://gandal.10.50.30.102.nip.io/} & Production \\
Documentation & \url{http://gandal.10.50.30.102.nip.io/docs} & Production \\
Port hôte direct & \url{http://10.50.30.102:8000/} & Diagnostic \\
\bottomrule
\end{tabularx}

\medskip
\noindent\textbf{Résolution DNS.} Le serveur s'appuie sur le service \emph{wildcard}
\textbf{nip.io} : toute entrée \code{<préfixe>.10.50.30.102.nip.io} résout
automatiquement vers \code{10.50.30.102}. Aucune configuration DNS supplémentaire n'est
donc nécessaire. En cas de bascule ultérieure sur un nom de domaine réel, il suffira de
pointer \code{gandal.<domaine>} (ou un enregistrement \emph{wildcard}) vers
\code{10.50.30.102}, sans modification côté serveur.

\medskip
\noindent\textbf{Compte d'administration} : utilisateur \code{admin} (rôle
\code{SuperAdmin}) -- mot de passe transmis séparément.

% =====================================================================
\section{Exploitation et maintenance}

Toutes les opérations s'effectuent depuis \code{/gandal-dev} :

\begin{lstlisting}[language=bash]
ssh root@10.50.30.102
cd /gandal-dev

# Etat des services
docker compose -f docker-compose.prod.yml ps

# Journaux applicatifs (suivi)
docker compose -f docker-compose.prod.yml logs -f api

# Mise a jour (apres synchronisation d'une nouvelle version du code)
docker compose -f docker-compose.prod.yml up -d --build

# Sauvegarde de la base de donnees
docker exec gandal-db pg_dump -U dc dc_backend > backup_$(date +%F).sql
\end{lstlisting}

\noindent Un script de déploiement automatisé est fourni : \code{scripts/deploy.sh}
(synchronisation \code{rsync} + reconstruction de l'image, avec gestion des coupures
réseau).

% =====================================================================
\section{Incidents rencontrés et résolutions}

\begin{longtable}{@{}p{0.30\linewidth} p{0.33\linewidth} p{0.30\linewidth}@{}}
\toprule
\textbf{Incident} & \textbf{Cause} & \textbf{Résolution} \\
\midrule
\endhead
Coupures SSH intermittentes & Instabilité du lien réseau vers l'hôte & Encapsulation des
commandes \code{ssh}/\code{scp} dans des boucles de réessai \\
Échec de \code{docker pull} & Résolveur DNS du serveur (\code{10.50.30.1:53})
intermittent & Réutilisation de l'image \code{postgres:17-alpine} déjà présente ; build
relancé en boucle \\
API non démarrée au 1\textsuperscript{er} essai & Base pas encore \code{healthy} & La
boucle de réessai a relancé \code{up} ; succès au 2\textsuperscript{e} essai \\
\code{register-admin} rejeté (422) & Domaine e-mail \code{.local} réservé & Adresse valide
\code{admin@enspy-gi.cm} \\
Build long interrompu & Lien réseau instable & Build exécuté \textbf{détaché}
(\code{nohup}), survivant aux ruptures de session \\
\bottomrule
\end{longtable}

% =====================================================================
\section{Risques résiduels et recommandations}

\begin{longtable}{@{}c >{\raggedright\arraybackslash}p{0.84\linewidth}@{}}
\toprule
\textbf{Priorité} & \textbf{Recommandation} \\
\midrule
\endhead
\textcolor{danger}{\textbf{Haute}} & \textbf{Activer HTTPS/TLS} sur le gateway
(certificat Let's Encrypt) -- le trafic transite actuellement en clair sur le port 80. \\
\textcolor{warn}{\textbf{Moyenne}} & \textbf{Restreindre ou fermer le port 8000 direct}
une fois la phase de diagnostic terminée. \\
\textcolor{warn}{\textbf{Moyenne}} & \textbf{Affiner \code{CORS\_ORIGINS}} avec l'origine
réelle du frontend de production. \\
\textcolor{warn}{\textbf{Moyenne}} & \textbf{Sauvegarde automatisée} du volume
\code{gandal\_pgdata} (cron \code{pg\_dump} + externalisation). \\
\textcolor{ok}{\textbf{Basse}} & \textbf{Surveiller l'espace disque} (87 \% utilisés) --
\code{docker system prune} régulier. \\
\textcolor{ok}{\textbf{Basse}} & \textbf{Stabiliser la résolution DNS} du serveur pour
fiabiliser les futurs \code{docker pull}. \\
\textcolor{ok}{\textbf{Basse}} & \textbf{Configurer le backend e-mail SMTP} si l'envoi de
notifications réelles est requis. \\
\textcolor{ok}{\textbf{Basse}} & \textbf{Activer Proxmox en production} : l'intégration est
désormais raccordée et validée contre le cluster réel (cf. \S\,10) ; il reste à ouvrir la
route réseau prod $\rightarrow$ hyperviseur puis à basculer \code{PROXMOX\_ENABLED=true}. \\
\bottomrule
\end{longtable}

% =====================================================================
\section{Annexes}

\textbf{Fichiers livrés sur le serveur} (\code{/gandal-dev/})
\begin{itemize}
  \item Code applicatif (synchronisé depuis le dépôt) ;
  \item \code{docker-compose.prod.yml} -- définition de la \emph{stack} de production ;
  \item \code{.env} -- variables d'environnement et secrets (non versionné).
\end{itemize}

\noindent\textbf{Fichier modifié sur le serveur}
\begin{itemize}
  \item \code{/root/iwm/ops/nginx/gateway.conf} -- ajout de la route \code{gandal.}
        (sauvegarde \code{.bak.<timestamp>} conservée).
\end{itemize}

\noindent\textbf{Références}
\begin{itemize}
  \item Dépôt : \url{https://github.com/Delmat237/GANDAL-API} ;
  \item Documentation API (Swagger) : \code{/docs} -- Redoc : \code{/redoc} ;
  \item Points de santé : \code{/health}, \code{/ready}.
\end{itemize}

\vspace{1.5cm}
\begin{center}
{\color{rule}\rule{0.5\linewidth}{0.6pt}}\\[0.3cm]
\textit{Rapport établi le 11 juin 2026 -- déploiement réalisé et validé.}
\end{center}

\end{document}
